% Source pp. 38-48 (Lectures1-6.pdf).
\section{Path Integrals}

Let me briefly review the method of path integral in ordinary quantum
mechanics. Here I will combine it with the idea of analytic continuation in
the time variable. We will start with the ``time evolution'' in pure imaginary
values of the time parameter. (Besides technical advantage of improving
convergence of the integral, the imaginary-time path integral makes evident
remarkable relation to classical statistical mechanics, see below.) The
real-time transition amplitudes can be recovered by analyticity. For
derivation directly in real time see \textbf{Sect.~9.1 of PS}.

\subsection*{Path integral in Quantum Mechanics}

Given a quantum mechanical system, with the space of states $\Ham$ and the
time-independent Hamiltonian $\hat H$, many problems can be reduced to the
calculations of appropriate matrix elements of its time-evolution operator,
\[
\bra{f}e^{-\frac i\hbar\hat Ht}\ket{i}\,.
\]
In terms of the stationary states $\ket n$ this can be written as
\[
\sum_n\braket{f}{n}\braket{n}{i}\,e^{-\frac i\hbar E_nt}\,.
\]
In a typical problem, the energy spectrum $E_n$ is bounded from below,
$E_n\ge E_0$, and $E_n$ grow to infinity with $n$. The above expression is a
sum of oscillating terms, and we would like to have an efficient mathematical
tool to sort them out. It is useful to consider this sum at complex values of
$t$. If the number of states does not grow exponentially with $E$ (and in
ordinary quantum system it never does), this sum converges absolutely if $t$
has negative imaginary part; thus it defines an analytic function of $t$ in
the lower half-plane. The real-time matrix element is the limiting value of
this analytic function. One can start with pure imaginary $t$,
\[
t=-i\tau\,,
\]
and try to evaluate the matrix elements of the corresponding bounded Hermitian
operator
\[
\bra{f}e^{-\frac1\hbar\hat H\tau}\ket{i}\,,
\]
anticipating that the real-time answer will be then obtained by analytic
continuation.

Consider a system with one degree of freedom $q$, and assume that the
Hamiltonian operator is ``local'' in the basis which diagonalizes $\hat q$;
this means that the matrix elements
\[
\bra{q}\hat H\ket{q'}
\]
have a support restricted to $q=q'$. Most typical Hamiltonians belong to this
class, for instance for $H=\frac12p^2+V(q)$ the above matrix element is
($\hbar=1$ from now on)
\[
-\frac12\,\delta''(q-q')+V(q)\,\delta(q-q')\,.
\]
In such cases it is possible to develop a path-integral representation for the
matrix element
\[
\bra{q_f}e^{-\hat H\tau}\ket{q_i}\,.
\]
The composition property (which follows from $\mathbb{I}=\int \dd q\,\ket q\bra q$)
\[
\bra{q_f}e^{-\hat H(\tau_1+\tau_2)}\ket{q_i}=\int\bra{q_f}e^{-\hat H\tau_1}\ket{q}\,\bra{q}e^{-\hat H\tau_2}\ket{q_i}\dd q
\]
can be applied repeatedly to reduce the above matrix element to the
composition of the matrix elements with arbitrarily small ``time'' intervals
$\Delta\tau$,
\[
\bra{q_f}e^{-\hat H\tau}\ket{q_i}=\int\prod_{k=1}^{n-1}\dd q_k\,\prod_{k=1}^n\bra{q_k}e^{-\hat H\Delta\tau}\ket{q_{k-1}}\,,
\]
where $q_0=q_i$, $q_n=q_f$, and $\Delta\tau=\tau/n$. In the limit of small
$\Delta\tau$, as the consequence of the ``locality'' of $\hat H$ in the
$q$-space, only the matrix elements $\bra{q'}\exp(-\hat H\Delta\tau)\ket{q}$
with $q'$ close to $q$ will bring significant contributions to the integral.

To see this, consider again the Hamiltonian of the form
\[
\hat H=\frac12\,\hat p^2+V(\hat q)\,.
\]
It is straightforward to show that in this case the matrix element
\[
G(q,q'|\tau)=\bra{q}e^{-\hat H\tau}\ket{q'}
\]
satisfies the differential equation of the diffusion type
\[
-\frac{\partial}{\partial\tau}\,G(q,q'|\tau)=\left(-\frac12\,\frac{\partial^2}{\partial q^2}+V(q)\right)\,G(q,q'|\tau)\,,
\]
with the initial condition
\[
G(q,q'|0)=\delta(q-q')\,.
\]
For $V(q)=0$ the problem is easily solved,
\[
G(q,q'|\tau)=\frac1{\sqrt{2\pi\tau}}\,e^{-\frac{(q-q')^2}{2\tau}}\,.
\]
This function indeed decays fast at $|q-q'|\gg\sqrt\tau$. In the case of
nonzero potential $V$ the problem still can be solved in the limit of small
$\tau$. We write
\[
G(q,q'|\tau)=e^{-\sigma(q,q'|\tau)}
\]
and look for $\sigma$ of the form of the following small-$\tau$ expansion
\[
\sigma(q,q'|\tau)=\frac{\sigma_0(q,q')}{\tau}+a\,\log\tau+b+\tau\,\sigma_1(q,q')+\dots\,.
\]
Substituting this ansatz into the above diffusion equation one finds (I
suggest doing it as an \textbf{Exercise})
\[
\sigma_0(q,q')=\frac{(q-q')^2}{2}\,,\qquad a=1/2\,,\qquad
\sigma_1(q,q')=\frac1{q-q'}\int_{q'}^qV(y)\dd y\,,\qquad\text{etc.}
\]
The leading terms do not depend on $V$, and are the same as for the free case,
and $\sigma_1$ is the potential $V$ averaged over the interval $[q',q]$. The
constant $b=\frac12\log2\pi$ is determined by the initial condition at
$\tau=0$.

Because of the factor $\exp(-\sigma_0(q,q')/\tau)$ the function
$G(q,q'|\tau)$ at small $\tau$ is sharply peaked around $|q-q'|=0$ with the
width $\sqrt\tau$. This is typical for the \emph{Brownian motion}, where
$\bar R\sim\sqrt t$ (our equation \emph{is} the diffusion equation). If the
potential $V$ is differentiable (which we assume), we have
\[
\sigma_1(q,q')=\frac12\,V(q)+\frac12\,V(q')+O(\tau)
\]
(I have used the fact that $(q-q')^2\sim\tau$). Therefore, for finite $\tau$
but large $n$ (i.e.\ small $\Delta\tau=\tau/n$) we can write
\[
\bra{q_f}e^{-\hat H\tau}\ket{q_i}=\lim_{n\to\infty}\left(\frac1{2\pi\Delta\tau}\right)^{n/2}\int\left(\prod_{k=1}^{n-1}\dd q_k\right)\,e^{-\mathcal A\{q_k\}}\,,
\]
where
\[
\mathcal A\{q_k\}=\sum_{k=1}^n\left[\frac{(q_k-q_{k-1})^2}{2\,\Delta\tau}+\Delta\tau\,\frac{V(q_k)+V(q_{k-1})}{2}\right]\,,
\]
and $\Delta\tau=\tau/n$, and it is understood that
\[
q_0=q_i\,,\qquad\text{and}\qquad q_n=q_f\,.
\]
Before the limit $n\to\infty$ is taken, the integration in this formula can be
understood as going over the piecewise linear ``trajectories'' $q(\tau)$
running from $q_i$ to $q_f$ (figure~\ref{fig:piecewise}).
\begin{figure}[ht]
\centering\figfile{piecewise}
\caption{A piecewise linear trajectory from $q_i$ at $\tau_0$ to $q_f$ at
$\tau_n$. The intermediate values $q_k=q(\tau_k)$ are integrated over.}
\label{fig:piecewise}
\end{figure}

When $n\to\infty$, $\Delta\tau\to0$, the function $\mathcal A\{q_k\}$ in the
exponential has formal limit
\[
\mathcal A[q(\tau)]=\int_0^\tau \dd\tau\,\left[\frac12\,\bigl(\dd q(\tau)/\dd\tau\bigr)^2+V\bigl(q(\tau)\bigr)\right]\,.
\]
It coincides with the ``imaginary time action'', i.e.\ the action associated
with our system, with $t$ replaced by $-i\tau$,
\[
e^{\frac i\hbar S[q(t)]}\ \xrightarrow{\ t\to-i\tau\ }\ e^{-\mathcal A[q(\tau)]}\,.
\]
It is conventional to write the $n\to\infty$ limit of the above $n$-fold
integral as
\[
\bra{q_f}e^{-\hat H(\tau_f-\tau_i)}\ket{q_i}=\int\bigl[Dq(\tau)\bigr]_{\substack{q(\tau_i)=q_i\\q(\tau_f)=q_f}}\,e^{-\mathcal A[q(\tau)]}\,,
\]
and to call it the (imaginary time) path integral, since the r.h.s.\ involves
the continuous path $q(\tau)$ from $q_i$ to $q_f$.

Few remarks.
\begin{enumerate}
\item The path integral method makes the most explicit use of the
\emph{superposition principle} of quantum theory. The latter states that if
some process can go several distinct ways, the full transition amplitude is
the sum of the amplitudes associated with each possible way. Here (when
returning back to real time) we sum the amplitudes associated with all
distinct paths from $q_i$ to $q_f$.

\item What kind of paths enter the path integral in the limit $n\to\infty$?
Although it is correct to think of the limiting paths as continuous functions
$q(\tau)$, absolute majority of the paths entering the integral are not smooth
curves. This can be observed directly from the fact that for
$\tau_1-\tau_2\to0$ we have $|q(\tau_1)-q(\tau_2)|\sim\sqrt{|\tau_1-\tau_2|}$,
not $\sim|\tau_1-\tau_2|$ as should be true for a continuously differentiable
function. From this point of view the two terms in the action play
significantly different role. The ``kinetic'' term $\frac12\int(\dd q/\dd\tau)^2\dd\tau$
selects the class of the paths entering the integral, i.e.\ those paths for
which $\bigl(q(\tau+\Delta\tau)-q(\tau)\bigr)^2/\Delta\tau$ remains finite as
$\Delta\tau\to0$; in this sense the factor
$\exp\bigl(-\frac12\int(\dd q/\dd\tau)^2\dd\tau\bigr)$ should be regarded as a part
of the functional measure. The factor $\exp\bigl(-\int V(q(\tau))\dd\tau\bigr)$
weights the paths according to their average potential energy.

\item Note that the symbol
\[
\bigl[Dq(\tau)\bigr]=\lim_{n\to\infty}\left(\frac1{2\pi\Delta\tau}\right)^{n/2}\prod_{k=1}^{n-1}\dd q_k
\]
contains divergent (in the limit) factor
\[
\left(\frac1{2\pi\Delta\tau}\right)^{n/2}=e^{-\frac{\log(2\pi\Delta\tau)}{2\,\Delta\tau}\,L}\,,
\]
where $L=(\tau_f-\tau_i)=n\,\Delta\tau$ is the ``volume'' of the imaginary
time interval. This factor is reminiscent of the irritating infinity we had
for the vacuum energy in the KG theory, and absorbing it into $[Dq(\tau)]$ is
analogous to subtracting $E_0$ from $\hat H$.

\item We have derived the above path integral representation of the evolution
operator in the imaginary time $\tau$. But one can repeat the above
calculation step by step, considering the paths in the real time $t$ instead
of $\tau$ (this is the form in which the path integral is usually introduced
in quantum mechanics, see for instance \S9.1 of the textbook). In this way one
arrives at the expression
\[
\bra{q_f}e^{-\frac i\hbar HT}\ket{q_i}=\int\bigl[Dq(t)\bigr]\,e^{\frac i\hbar S[q(t)]}\,,
\]
with the phase $e^{\frac i\hbar S}$ replacing $e^{-\mathcal A}$. This of
course is equivalent approach, although, unlike the imaginary time scheme, the
integrals involved are not absolutely convergent. Usually, the most efficient
way to handle such integrals is through analytic continuation (in variables or
parameters) which makes the integral absolutely convergent, with subsequent
taking the limiting values. The imaginary-time approach simply makes this
procedure explicit.

\item Generalization of the above construction to the systems with many degrees
of freedom $q_a$ is straightforward, if the Hamiltonian is of the form
\[
H=\sum_a\frac12\,p_a^2+V(q_a)\,;
\]
one obtains
\[
\int\bigl[Dq_a\bigr]\,e^{-\mathcal A[q_a(\tau)]}
\]
for the evolution operator. In more complicated cases the path integral in the
phase space can be useful; this is briefly discussed in Sect.~9.1 of PS.
\end{enumerate}

\subsection*{Relation to Classical Statistical Mechanics}

At this point, let me deviate to discuss remarkable relation which exists
between the \textbf{Path Integral in Quantum Mechanics} (in its imaginary time
version) and the \textbf{Classical Statistical Mechanics}. Let us recall some
generalities of the classical statistical mechanics.

Suppose we have a classical dynamical system with the phase-space coordinates
$\{P_i,Q_i\}$ (I will use capitals here because the relation between these
$Q_i$ and $q_a$ appearing in the path integral is not going to be
straightforward), and the Hamiltonian $\Ham(P_i,Q_i)$. In a thermal
equilibrium state at a temperature $T$ the probability distribution of
microscopic states is given by the Gibbs formula
\[
\mathcal P(P_i,Q_i)\prod_i \dd P_i\dd Q_i=Z^{-1}\,e^{-\beta\Ham(P_i,Q_i)}\,\frac1{N!}\prod_i \dd P_i\dd Q_i\,,
\]
where $\beta=1/kT$, and $Z$ is the partition function
\[
Z=\int\frac{\prod_i \dd P_i\dd Q_i}{N!}\,e^{-\beta\Ham(P_i,Q_i)}\,.
\]
Calculating the partition function is one of the main problems in statistical
mechanics because the thermodynamic quantities are expressed through $Z$, say
the free energy is $\beta F=-\log Z$. Again, a typical form of the Hamiltonian
is
\[
\Ham(P_i,Q_i)=\sum_i\frac12\,P_i^2+U(\{Q_i\})\,.
\]
In this case the integration over $P_i$ is easy to perform,
\[
Z=\frac1{N!}\left(\frac{2\pi}{\beta}\right)^{\frac N2}Z_{\text{conf}}\,,
\]
where $N$ is the number of the degrees of freedom, $i=1,2,\dots,N$, and
$Z_{\text{conf}}$ is the configuration-space integral
\[
Z_{\text{conf}}=\int\prod_i \dd Q_i\,e^{-\beta U(\{Q_i\})}
\]
which contains most of interesting physics.

As a special example of such system, consider an elastic string whose
configuration is described by a function $q(\tau)$, $\tau\in[0,L]$; the values
$q(\tau)$ for all $\tau$ within this interval play the role of the
configuration-space coordinates (figure~\ref{fig:string}),
\[
\{Q_i\}\to q(\tau)\,.
\]
\begin{figure}[ht]
\centering\figfile{string}
\caption{A configuration $q(\tau)$ of the elastic string on $[0,L]$ with the
ends held at $q(0)=q_i$ and $q(L)=q_f$.}
\label{fig:string}
\end{figure}
We also assume that the string lies in a potential well $V(q)$, so that the
potential energy of the string is
\[
U[q(\tau)]=\int_0^L\dd\tau\,\left[V\bigl(q(\tau)\bigr)+\frac12\,\left(\frac{\dd q(\tau)}{\dd\tau}\right)^2\right]\,,
\]
where the last term accounts for the elastic energy; to simplify the equations
I have assumed that the elastic string tension equals 1.

To evaluate the configuration part of the partition function, we must
integrate over the functions $q(\tau)$ satisfying given boundary conditions at
the ends of the interval $[0,L]$, say
\[
q(0)=q_i\,,\qquad q(L)=q_f\,,
\]
with the statistical weights $\exp(-\beta\,U[q(\tau)])$. Thus
\[
Z_{\text{conf}}=\int\bigl[Dq(\tau)\bigr]\,e^{-\beta U[q(\tau)]}\,.
\]
With the above form of $U[q(\tau)]=\mathcal A[q(\tau)]$, this expression is
identical to the path integral for the imaginary time transition amplitude,
\[
Z_{\text{conf}}\sim\bra{q_f}e^{-\frac1\hbar\hat HL}\ket{q_i}\,,
\]
provided we also identify
\[
\beta=1/\hbar\,.
\]
We see that quantum mechanics with a single degree of freedom $q$, and the
Hamiltonian operator
\[
\hat H=\frac12\,\hat p^2+V(\hat q)\,,
\]
is related in this way to classical statistical mechanics of a system with
continuously many degrees of freedom $q(\tau)$.

\subsection*{Correlation Functions}

In statistical mechanics one may also be interested in \emph{correlation
functions}. In our example one may ask how the string positions $q(\tau)$ at
different points, say at $\tau_1,\tau_2,\dots,\tau_N$, are correlated in the
thermal equilibrium state. In figure~\ref{fig:stringcorr} I made a trivial
shift of the variable $\tau$, so that from now on $\tau\in[-L/2,L/2]$.
\begin{figure}[ht]
\centering\figfile{stringcorr}
\caption{The string on $[-L/2,L/2]$, with its positions $q(\tau_1),\dots,q(\tau_N)$
marked at $\tau_1\le\tau_2\le\dots\le\tau_N$.}
\label{fig:stringcorr}
\end{figure}

The correlations are described by the averaged products
$q(\tau_1)q(\tau_2)\cdots q(\tau_N)$, which in turn are given by the integrals
\[
\bigl\langle q(\tau_1)\,q(\tau_2)\cdots q(\tau_N)\bigr\rangle
=Z^{-1}\int q(\tau_1)q(\tau_2)\cdots q(\tau_N)\,e^{-\mathcal A[q(\tau)]}\,\bigl[Dq(\tau)\bigr]\,,
\]
because $Z^{-1}e^{-\mathcal A[q(\tau)]}\,[Dq(\tau)]$ is the statistical
weight of the microstate $q(\tau)$. It is instructive to find out how these
correlation functions are interpreted in associated quantum mechanical
problem.

Without loss of generality we can assume that
\[
-L/2\le\tau_1\le\tau_2\le\dots\le\tau_N\le L/2\,.
\]
The integration over all microstates $q(\tau)$ can be performed in two steps:
First, we fix the values of $q(\tau)$ at the points $\tau_1,\tau_2,\dots,\tau_N$
to be
\[
q(\tau_1)=q_1\,,\quad q(\tau_2)=q_2\,,\quad\dots\quad q(\tau_N)=q_N\,,
\]
and integrate over all $q(\tau)$ subject to these constraints. After that, the
integration over $q_1,q_2,\dots,q_N$ is performed to complete the integral.
Comparing this with the path integral representation for the evolution
operator, say
\[
\bra{q_2}e^{-\hat H(\tau_2-\tau_1)}\ket{q_1}=\int\bigl[Dq(\tau)\bigr]_{\substack{q(\tau_1)=q_1\\q(\tau_2)=q_2}}\,e^{-\mathcal A[q(\tau)]}\,,
\]
we find
\begin{multline*}
\bigl\langle q(\tau_1)q(\tau_2)\cdots q(\tau_N)\bigr\rangle=Z^{-1}\int\left[\prod_{k=1}^N\dd q_k\right]\times q_1q_2\cdots q_N\\
\times\bra{q_f}e^{-\hat H(L/2-\tau_N)}\ket{q_N}\bra{q_N}e^{-\hat H(\tau_N-\tau_{N-1})}\ket{q_{N-1}}\cdots\bra{q_1}e^{-\hat H(\tau_1+L/2)}\ket{q_i}\,.
\end{multline*}
Recall that in the coordinate representation the operator $\hat q$ acts by
multiplication
\[
\hat q\,\Psi(q)=q\,\Psi(q)\,,
\]
so that in this expression we can replace
\[
\int \dd q\,\ket{q}\,q\,\bra{q}=\hat q\,.
\]
Then the above expression can be written as
\[
\bra{q_f}e^{-\hat H(L/2-\tau_N)}\,\hat q\,e^{-\hat H(\tau_N-\tau_{N-1})}\,\hat q\cdots\hat q\,e^{-\hat H(\tau_1+L/2)}\ket{q_i}
\]
which is best expressed in terms of the ``Euclidean'' (i.e.\ imaginary time)
Heisenberg operators
\[
\hat q_E(\tau)=e^{\hat H\tau}\,\hat q\,e^{-\hat H\tau}
\]
($\tau=it$), allowing one to write the above formula as
\[
\bra{q_f}e^{-\hat H\frac L2}\,\hat q_E(\tau_N)\,\hat q_E(\tau_{N-1})\cdots\hat q_E(\tau_1)\,e^{-\hat H\frac L2}\ket{q_i}\,.
\]
To simplify things further, let's take the limit $L\to\infty$, thus turning to
the correlation functions of the infinite string. In this limit the exact
boundary conditions, e.g.\ the values $q_i$ and $q_f$, become unimportant.
Indeed, the states $\ket{q_i}$ and $\ket{q_f}$ can be expanded in the
stationary states $\ket n$ of the Hamiltonian $\hat H$, say
\[
\ket{q_i}=\sum_n\ket n\braket{n}{q_i}=\sum_n\ket n\,\Psi_n^*(q_i)\,,
\]
where it is assumed that
\[
\braket{n}{m}=\delta_{n,m}\,.
\]
In the limit $L\to\infty$ the contribution of the ground state $\ket0$
dominates, since all other states are suppressed by the factor
$\sim e^{-(E_n-E_0)\frac L2}$. Thus we have
\[
\bigl\langle q(\tau_1)q(\tau_2)\cdots q(\tau_N)\bigr\rangle=\frac{\braket{q_f}{0}\braket{0}{q_i}}{Z}\,e^{-E_0L}\,
\bra0\hat q_E(\tau_N)\,\hat q_E(\tau_{N-1})\cdots\hat q_E(\tau_1)\ket0\,.
\]
Furthermore, note that
\[
Z=\bra{q_f}e^{-\hat HL}\ket{q_i}\ \to\ \braket{q_f}{0}\braket{0}{q_i}\,e^{-E_0L}\,,
\]
hence the first factor in the above equation equals 1 (in the limit
$L=\infty$), and we have
\[
\bigl\langle q(\tau_1)\cdots q(\tau_N)\bigr\rangle=\bra0\hat q_E(\tau_N)\,\hat q_E(\tau_{N-1})\cdots\hat q_E(\tau_1)\ket0\,.
\]
Note that the operators here are arranged in the order of $\tau$ increasing
from the right to the left. Let me stress that the Euclidean Heisenberg
operators can not be arranged any other way, as otherwise the product of the
operators (the sum over the intermediate states) does not converge. More
generally, if we consider complex values of $\tau_i$ (which we will), the above
product of the operators makes direct sense in the domain
\[
\Re\tau_1\le\Re\tau_2\le\dots\le\Re\tau_N
\]
where the sums over the intermediate states converge (hence the sums yield
analytic function).

To relate the above Euclidean correlation function to the real-time ones one
analytically continues to pure imaginary values of $\tau$,
\[
\tau_k=it_k+0\,,
\]
where ``$+0$'' shift indicates that if there is a branch cut along the
imaginary axis the values at the right edge are to be taken.
